% Source PDF: source/普通高考/2014/2014江苏.pdf, page 1.
% Node -> directed-edge list: start -> n=0 -> n=n+1 -> 2^n>20?;
% 2^n>20? --Y--> output n -> finish; --N--> right loop-back -> main stem above n=n+1.
% The loop return has a single left-pointing arrow into the main stem; the
% independent main-stem arrow continues downward into the update box.
\begin{tikzpicture}[
  >=Stealth,
  line width=0.45pt,
  line cap=butt,
  line join=miter,
  every node/.style={font=\normalsize,anchor=center,align=center,
    execute at begin node={\setlength{\parindent}{0pt}}},
  flow/.style={draw=black,anchor=center,align=center,inner sep=0pt,
    execute at begin node={\setlength{\parindent}{0pt}}},
]
  \node[flow,rounded corners=0.18cm,minimum width=1.50cm,minimum height=0.72cm]
    (start) at (0,10.6) {开始};
  \node[flow,rectangle,minimum width=1.95cm,minimum height=0.80cm]
    (init) at (0,9.10) {$n\leftarrow 0$};
  \node[flow,rectangle,minimum width=2.95cm,minimum height=0.86cm]
    (increment) at (0,7.35) {$n\leftarrow n+1$};
  \node[flow,diamond,aspect=2.75,minimum width=4.20cm,minimum height=1.10cm]
    (test) at (0,5.55) {$2^n>20$};
  \path[flow] (-1.22,3.88) -- (1.22,3.88) -- (0.96,3.08)
    -- (-1.48,3.08) -- cycle;
  \node[font=\normalsize] (output) at (-0.20,3.48) {输出 $n$};
  \node[flow,rounded corners=0.18cm,minimum width=1.50cm,minimum height=0.72cm]
    (finish) at (0,1.95) {结束};

  \draw[->] (start.south) -- (init.north);
  \draw[->] (init.south) -- (increment.north);
  \draw[->] (increment.south) -- (test.north);
  \draw[->] (test.south) -- node[right=2pt] {Y} (0,3.88);
  \draw[->] (0,3.08) -- (finish.north);

  % N branch loops right and returns with one horizontal arrow into the main stem.
  \draw (test.east) -- (2.80,5.55) -- (2.80,8.25);
  \draw[->] (2.80,8.25) -- (0,8.25);
  \node[anchor=west,inner sep=0pt] at (1.95,5.88) {N};
\end{tikzpicture}
